\subsection{COMPARISON OF THEORIES}

A theory \(\mathcal{C}'\) is said to be stronger than a theory \(\mathcal{C}\) if all the signs of \(\mathcal{C}\) are signs of \(\mathcal{C}'\) , all the explicit axioms of \(\mathcal{C}\) are theorems in \(\mathcal{C}'\) , and the schemes of \(\mathcal{C}\) are schemes of \(\mathcal{C}'\) .

C4. If a theory \(\mathfrak{G}'\) is stronger than a theory \(\mathfrak{G}\) , then all the theorems of \(\mathfrak{G}\) are theorems of \(\mathfrak{G}'\) .

Let \(R_1, R_2, \ldots, R_n\) be a proof in \(\mathfrak{G}\) . We shall show, step by step, that each \(R_i\) is a theorem in \(\mathfrak{G}'\) . Suppose that this is true for the relations preceding \(R_k\) . If \(R_k\) is an axiom of \(\mathfrak{G}\) , it is a theorem in \(\mathfrak{G}'\) by hypothesis. If \(R_k\) is preceded by relations \(R_i\) and \(R_i \Longrightarrow R_k\) , we know already that \(R_i\) and \(R_i \Longrightarrow R_k\) are theorems in \(\mathfrak{G}'\) , and therefore \(R_k\) is a theorem in \(\mathfrak{G}'\) by virtue of C1. Hence, in every case, \(R_k\) is a theorem in \(\mathfrak{G}'\) , and the proof is complete.

If each of two theories \(\mathfrak{G}\) and \(\mathfrak{G}'\) is stronger than the other, \(\mathfrak{G}\) and \(\mathfrak{G}'\) are said to be equivalent. Then every theorem of \(\mathfrak{G}\) is a theorem of \(\mathfrak{G}'\) , and vice versa.

C5. Let \(\mathfrak{G}\) be a theory, let \(A_{1}, A_{2}, \ldots, A_{n}\) be its explicit axioms, \(a_{1}, a_{2}, \ldots, a_{h}\) its constants, and let \(T_{1}, T_{2}, \ldots, T_{h}\) be terms in \(\mathfrak{G}\) . Suppose that

\[
\left(\boldsymbol {T} _ {1} \mid \boldsymbol {a} _ {1}\right) \left(\boldsymbol {T} _ {2} \mid \boldsymbol {a} _ {2}\right) \dots \left(\boldsymbol {T} _ {h} \mid \boldsymbol {a} _ {h}\right) \boldsymbol {A} _ {i} \quad (\text { for } i = 1, 2, \dots , n)
\]

are theorems in a theory \(\mathfrak{G}'\) , that the signs of \(\mathfrak{G}\) are signs of \(\mathfrak{G}'\) , and that the schemes of \(\mathfrak{G}\) are schemes of \(\mathfrak{G}'\) . Then, if \(A\) is a theorem in \(\mathfrak{G}\) ,

\[
\left(\boldsymbol {T} _ {1} \mid \boldsymbol {a} _ {1}\right) \dots \left(\boldsymbol {T} _ {h} \mid \boldsymbol {a} _ {h}\right) \boldsymbol {A}
\]

is a theorem in \(\mathfrak{C}'\) .

For \(\mathfrak{c}'\) is stronger than the theory \((T_1|a_1) \ldots (T_h|a_h)\mathfrak{c}\) , and we can apply C2 and C4.

When we use this procedure to deduce a theorem in \(\mathfrak{C}'\) from a theorem in \(\mathfrak{C}\) , we say that we are applying in \(\mathfrak{C}'\) the results of \(\mathfrak{C}\) . Intuitively, the axioms of \(\mathfrak{C}\) express properties of \(a_1, a_2, \ldots, a_h\) , and \(A\) expresses a property which is a consequence of these axioms. If the objects \(T_1, T_2, \ldots, T_h\) in \(\mathfrak{C}'\) have the properties expressed by the axioms of \(\mathfrak{C}\) , then they also have the property \(A\) .

*For example, in the theory of groups \(\mathfrak{C}\) , the explicit axioms contain two constants \(G\) and \(\mu\) (the group and the law of composition). In the theory of sets \(\mathfrak{C}'\) , we define two terms: the real line and addition of real numbers. If we substitute these terms for \(G\) and \(\mu\) respectively in the explicit axioms of \(\mathfrak{C}\) , we obtain theorems in \(\mathfrak{C}'\) . Moreover, the schemes and signs of \(\mathfrak{C}\) and \(\mathfrak{C}'\) are the same. We may therefore ``apply the results of group theory to the additive group of real numbers''. We say that we have constructed a model for group theory in the theory of sets. (Note that since the theory of groups is stronger than the theory of sets, we can also apply the results of the theory of sets to the theory of groups.)*

Remark. Under the hypotheses of C5, if the theory \(\mathfrak{C}\) turns out to be contradictory, the same will be true of \(\mathfrak{C}'\) . For if \(A\) and ``not \(A\)'' are theorems in \(\mathfrak{C}\) , then \((T_1|a_1)\ldots (T_h|a_h)A\) and \(\mathrm{not}(T_1|a_1)\ldots (T_h|a_h)A\) are theorems in \(\mathfrak{C}'\) . * For example, if the theory of groups were contradictory, the theory of sets would also be contradictory. *

